\chapter{Polycrisis and Metacrisis}
\label{ch:03}

The vocabulary of crisis has expanded in recent decades in a way that is partly descriptive and partly rhetorical, and a book that proposes to organize its argument around an inventory of problems must first say what it takes the inventory to be an inventory of. The word polycrisis, which entered wide use in the scholarly literature through the work of Morin and Kern~\cite{morin1999} and was later revived in discussions of global risk~\cite{lawrence2022}, names the situation in which several distinct crises, each of which might be manageable if it occurred alone, interact so that the whole exceeds the sum of its parts. The word metacrisis, which has a less settled meaning and is used by different authors for different purposes, is employed here in a narrow and deliberately modest sense: it names the condition in which the capacity of institutions to perceive, deliberate about, and correct problems is itself degraded by the problems it is called upon to address. The two terms are not synonyms, and the difference between them determines what kind of help can be expected from any increase in technical capability. A polycrisis is a fact about the structure of coupling among domains, and one can in principle map it, measure the strength of the couplings, and identify the places at which a modest intervention relieves many pressures at once. A metacrisis, as defined here, is a fact about the corrective apparatus, and it is the more worrying of the two because a failure of correction cannot be repaired by the instrument that has failed. The coupling of risks across domains has been studied in the literature on globally networked risks~\cite{helbing2013}, and its feedback structure in the tradition of systems dynamics~\cite{meadows2008}.

The coupling among domains is easy to illustrate and difficult to quantify. Drought reduces agricultural yield, reduced yield raises food prices, higher prices produce unrest in importing regions, unrest displaces populations, displacement places pressure on the water and health systems of the regions that receive them, and strained health systems are less able to contain the infectious diseases that spread more readily among displaced and malnourished persons. None of these links is mysterious, and each has a literature. What is distinctive is the closure of the loop, since any effect that returns to amplify its own cause raises the possibility that a system which is stable with respect to each of its components taken separately is unstable as a whole. The literature on global systemic risk, including the survey by Lawrence, Janzwood, and Homer-Dixon, emphasizes this point and also emphasizes a limitation: the couplings are poorly measured, the assumption that they are linear is a convenience, and the identification of cascades after the fact is much easier than their anticipation. A responsible treatment therefore resists two temptations. The first is to treat the interdependence as a reason for despair, as if the existence of loops implied that no intervention could succeed. The second is to treat it as a reason for a single grand solution, as if the existence of loops implied that one intervention at the right point would resolve all of them. Neither follows. What follows is that interventions must be evaluated with their feedback effects in view, and that the identification of the points at which a loop can be cut cheaply is among the most valuable analytical tasks that can be assigned to systems with high capability for modelling.

It is at this point that the relevance of the Gebru interview to the technical program becomes clear. A common move in the rhetoric of redemptive artificial intelligence is to name the crises, to assert that sufficiently capable systems will solve them, and to omit the mechanism. The omission is conspicuous in precisely the cases where the crises are coupled, because a claim to have solved one member of a coupled set must also account for the effect of the solution on the others, and a claim to have solved them all must account for the institutional process by which the solutions will be adopted, funded, and maintained. Capability in modelling does not remove the difficulty, but it does bear on it in two ways that can be stated without exaggeration. It can reduce the cost of constructing and exploring models of coupled systems, so that the sensitivity of an outcome to an intervention can be examined across many scenarios rather than a few. And it can assist in the search for interventions that exploit asymmetries in the coupling, for example by identifying a constraint whose relaxation lowers several coupling coefficients at once. It cannot determine which of the several defensible weightings of the affected goods ought to govern the choice, and it cannot supply the legitimacy without which a solution that is technically adequate will be resisted, evaded, or reversed. Because the latter are the usual causes of failure in practice, the book treats the modelling contribution as real but bounded and does not allow it to stand in for institutional work.

The metacrisis in the narrow sense adopted here has a more precise structure. Institutions correct problems by detecting them, by deliberating about responses, by mobilizing resources, and by implementing and monitoring the response. Each of these functions depends on resources that are themselves consumed or damaged by sustained stress. Attention is scarce, and a population under continuous emergency has less of it to give to long-term matters. Trust is a stock that is built slowly and depleted rapidly, and a series of failures of public bodies reduces the willingness of citizens to accept the costs of the next corrective measure. Fiscal capacity is reduced when the economy is damaged, and administrative competence is reduced when experienced personnel depart. Where these dependencies are present, the response to stress diminishes as stress grows, and the system possesses two qualitatively different regimes: one in which stress remains bounded because the corrective capacity is adequate, and another in which stress grows without bound because the capacity has been eroded below the level needed to arrest it. Between the regimes lies a threshold, and the characteristic feature of the threshold is that it is not visible from within the lower regime, because the indicators of stress remain moderate until the point at which the upper regime has already been entered. This is the sense in which the metacrisis is a distinct problem and not merely a large collection of ordinary ones.

The implications for the program of this book are of three kinds. The first concerns priorities. If corrective capacity can be eroded by the stress it is meant to relieve, then interventions that protect or restore that capacity, such as the maintenance of verification institutions, the preservation of public trust through honest reporting of failures, and the avoidance of dependence on single points of failure, have a value that is not captured by their direct effects. The second concerns the use of advanced systems. A system that is introduced into institutions in a way that reduces the capacity of human beings to understand and contest decisions may increase efficiency in the short run while lowering corrective capacity in the long run, and the architecture of authority described in Chapter 13 is in part a safeguard against that outcome. The third concerns the temptation to treat the existence of a crisis as a reason to suspend the safeguards. The argument that an emergency of sufficient gravity justifies the delegation of authority to a capable system, with accountability to be reintroduced later, is the form that the displacement of responsibility takes in the rhetoric of urgency, and it is precisely the argument the framework is constructed to resist. A program that can be defended only on the assumption that the emergency will excuse its omissions is not a program that has met its own standard.

A final clarification concerns the relation of the present usage to broader and more ambitious theories of civilizational crisis. Some authors use the term metacrisis for a crisis of meaning or of worldview, and others use it for a purported common root of all the particular crises. The present book takes no position on whether such common roots exist. It restricts attention to a structural property that can be stated, tested in models, and in principle observed, namely the dependence of corrective capacity on the state of the system being corrected. Where a larger interpretation can be built on that property, it is welcome, but the program does not depend on it, and a reader who rejects every larger interpretation is not thereby released from the practical conclusion that correction capacity is an asset to be protected.

\section*{Formalization}

Two models are developed, the first for the coupling that defines a polycrisis and the second for the erosion of corrective capacity that defines a metacrisis in the sense above. Both are deliberately minimal, and the purpose of each is to display a mechanism whose qualitative behavior survives enrichment of the model, not to predict.

Consider $n$ domains with stress vector $x_t\in\R_{\ge0}^n$ and a nonnegative coupling matrix $M=(m_{ij})$, where $m_{ij}$ is the stress induced in domain $i$ per unit of stress in domain $j$ in the next period. Let $e\in\R^n_{\ge0}$ denote exogenous stress per period and suppose the stress evolves as $x_{t+1}=Mx_t+e$ with $x_0=0$. The domains are \emph{separately stable} if the diagonal part $D=\mathrm{diag}(m_{11},\dots,m_{nn})$ satisfies $\rho(D)<1$, where $\rho$ denotes the spectral radius, and the system is \emph{jointly stable} if $\rho(M)<1$.

\begin{definition}[Polycrisis]
The system is in a \emph{polycrisis regime} if $\rho(D)<1$ and $\rho(M)\ge 1$, that is, if each domain would be stable if isolated from the others and the coupled system is not.
\end{definition}

\begin{proposition}
Let $M\ge0$ and $e\ge0$. (i) If $\rho(M)<1$ then $x_t\to x^\ast=(I-M)^{-1}e=\sum_{k\ge0}M^ke$, and $x^\ast\ge e$ componentwise. Moreover, if $\|M\|_1<1$, where $\|M\|_1$ is the maximum column sum, then $\mathbf 1^\top x^\ast\le \mathbf 1^\top e/(1-\|M\|_1)$. (ii) If $\rho(M)\ge1$ and $w\ge0$ is a left Perron vector of $M$ with $w^\top e>0$, then $\mathbf 1^\top x_t\to\infty$. (iii) For $n=2$ with $D=0$ and $M=\left(\begin{smallmatrix}0&a\\b&0\end{smallmatrix}\right)$, one has $\rho(M)=\sqrt{ab}$, so the coupled system is unstable exactly when $ab\ge1$ although $\rho(D)=0$.
\end{proposition}

\begin{proof}
For (i), $x_t=\sum_{k<t}M^ke$ by induction, and the Neumann series converges when $\rho(M)<1$ to $(I-M)^{-1}$, which is nonnegative because every term is. The term $k=0$ gives $x^\ast\ge e$. For the column-sum bound, $\|M^ke\|_1\le\|M\|_1^k\|e\|_1$ for nonnegative vectors, whence $\mathbf 1^\top x^\ast\le\sum_k\|M\|_1^k\,\mathbf 1^\top e$. For (ii), by the Perron--Frobenius theorem for nonnegative matrices, the left Perron vector $w\ge0$ satisfies $w^\top M=\rho(M)w^\top$, so $w^\top x_t=\sum_{k<t}\rho(M)^kw^\top e$, which diverges if $\rho(M)\ge1$ and $w^\top e>0$, and $\mathbf 1^\top x_t\ge w^\top x_t/\max_iw_i$. For (iii), the characteristic polynomial of $M$ is $\lambda^2-ab$.
\end{proof}

For example, with $a=0.8$ and $b=1.4$ each domain responds to the other with a gain of the order of one, but $ab=1.12$, so $\rho(M)=\sqrt{1.12}\approx1.058$ and the coupled system diverges, whereas with $a=0.5$ and $b=0.8$ one has $\rho(M)\approx0.632$ and the system settles at a finite amplification of the exogenous stress. The point of the example is that the stability of a polycrisis is not determined by any single coefficient and can be restored by reducing either, which is the formal content of the statement that cutting a loop at its cheapest point is a legitimate strategy.

For the metacrisis, let a scalar $x\ge0$ denote aggregate stress and let corrective capacity be $\kappa(x)=\kappa_0-\gamma x$ for $x<\kappa_0/\gamma$, with $\kappa_0,\gamma>0$, so that capacity is eroded linearly by stress. With exogenous inflow $e>0$ and self-amplification coefficient $a\ge0$ with $a<\kappa_0$, let
\[
\dot x=e+a x-\kappa(x)\,x=e-(\kappa_0-a)x+\gamma x^2 .
\]

\begin{proposition}
Define $e^\ast=(\kappa_0-a)^2/(4\gamma)$. If $e<e^\ast$, the equation has two equilibria $x_\pm=\bigl[(\kappa_0-a)\pm\sqrt{(\kappa_0-a)^2-4\gamma e}\bigr]/(2\gamma)$ with $0<x_-<x_+$; $x_-$ is asymptotically stable, $x_+$ is unstable, and every trajectory starting at $x_0<x_+$ converges to $x_-$ while every trajectory starting at $x_0>x_+$ diverges. If $e>e^\ast$, no equilibrium exists and every trajectory diverges. The width of the safe basin, $x_+-x_-=\sqrt{(\kappa_0-a)^2-4\gamma e}/\gamma$, decreases strictly to zero as $e\uparrow e^\ast$.
\end{proposition}

\begin{proof}
Equilibria solve $\gamma x^2-(\kappa_0-a)x+e=0$, whose discriminant is $(\kappa_0-a)^2-4\gamma e$, positive iff $e<e^\ast$. The right-hand side $f(x)=\gamma x^2-(\kappa_0-a)x+e$ is an upward parabola, negative between its roots and positive outside them, so $f'(x_-)<0<f'(x_+)$ and the sign of $f$ determines the stated convergence and divergence. If $e>e^\ast$ then $f>0$ everywhere and $x$ increases without bound.
\end{proof}

With $\kappa_0=1$, $a=0.3$, and $\gamma=0.5$ one finds $e^\ast=0.49/2=0.245$. At $e=0.1$ the equilibria are $x_-\approx0.161$ and $x_+\approx1.239$, so the safe basin is wide. At $e=0.2$ they are $x_-=0.4$ and $x_+=1.0$, and the stressed but stable state sits at $0.4$, only $0.6$ below the point beyond which recovery fails. The calculation exhibits the property that motivated the definition: the equilibrium stress rises gradually as $e$ increases, from $0.161$ to $0.4$ in this example, while the margin to the irreversible regime shrinks at a faster relative rate, so that observation of the equilibrium alone understates the proximity of the threshold. The implication for institutional design is that the proper quantity to monitor is the margin $x_+-x_-$, or its proxy in the form of the reserve of corrective capacity, and not the level of stress.
